% =========================================================================
% PRACTICE EXAM 1 (2027 EXAM READY DIAGNOSTIC)
% =========================================================================
\chapter{Practice Exam 1}

\section*{Exam Overview \& Instructions}
This authentic practice exam simulates the official AP Statistics Exam.
\begin{itemize}[leftmargin=1.2em]
  \item \textbf{Section I:} 40 Multiple-Choice Questions (90 Minutes | 50\% of Total Score).
  \item \textbf{Section II:} 6 Free-Response Questions (90 Minutes | 50\% of Total Score).
  \item Calculator active throughout both sections. Standard normal, $t$, and $\chi^2$ tables provided.
\end{itemize}

\section{Section I: Multiple-Choice Questions (Questions 1--40)}
\begin{enumerate}[leftmargin=1.5em]
  \item A set of 20 exam scores has $\text{Mean} = 75$ and $\text{SD} = 8$. A score of 91 corresponds to a $z$-score of:
  \begin{enumerate}[label=(\Alph*)]
    \item 1.50 \item 2.00 \item 2.25 \item 1.75 \item 1.25
  \end{enumerate}

  \item Which of the following is a categorical variable?
  \begin{enumerate}[label=(\Alph*)]
    \item Annual income \item Blood type (A, B, AB, O) \item Systolic blood pressure \item Commute time in minutes \item High school GPA
  \end{enumerate}

  \item A scatterplot exhibits a strong negative linear association with $r = -0.92$. The coefficient of determination $r^2$ is:
  \begin{enumerate}[label=(\Alph*)]
    \item $-0.846$ \item 0.846 \item $-0.92$ \item 0.92 \item 0.080
  \end{enumerate}

  \item In a modified boxplot, the whiskers extend to:
  \begin{enumerate}[label=(\Alph*)]
    \item The outlier boundary fences $Q_1 - 1.5(\text{IQR})$ and $Q_3 + 1.5(\text{IQR})$
    \item The absolute minimum and maximum data values
    \item The most extreme observations that fall within the boundary fences
    \item $Q_1$ and $Q_3$
    \item $\bar{x} \pm 2s$
  \end{enumerate}

  \item An observational study finds people who take multivitamins have fewer colds. Can we conclude multivitamins prevent colds?
  \begin{enumerate}[label=(\Alph*)]
    \item Yes, if $n > 1000$ \item No, because confounding variables cannot be controlled \item Yes, if $p < 0.05$ \item No, because vitamins are categorical \item Yes, if double-blind
  \end{enumerate}

  \item If $P(A) = 0.5, P(B) = 0.4,$ and $A$ and $B$ are independent, what is $P(A \cup B)$?
  \begin{enumerate}[label=(\Alph*)]
    \item 0.90 \item 0.70 \item 0.20 \item 0.50 \item 0.80
  \end{enumerate}

  \item A random variable $X \sim B(50, 0.20)$. The mean and standard deviation are:
  \begin{enumerate}[label=(\Alph*)]
    \item $\mu = 10, \sigma = 2.83$ \item $\mu = 10, \sigma = 8.00$ \item $\mu = 25, \sigma = 5.00$ \item $\mu = 10, \sigma = 2.00$ \item $\mu = 50, \sigma = 10$
  \end{enumerate}

  \item The Central Limit Theorem applies to which of the following?
  \begin{enumerate}[label=(\Alph*)]
    \item The shape of the population \item The shape of a single sample \item The sampling distribution of sample means $\bar{x}$ as $n \ge 30$ \item The sampling distribution of $\hat{p}$ \item The variance of a binomial variable
  \end{enumerate}

  \item In testing $H_0: p = 0.5$ vs. $H_a: p > 0.5$, a sample gives $z = 2.33$. The $p$-value is:
  \begin{enumerate}[label=(\Alph*)]
    \item 0.0198 \item 0.0099 \item 0.9901 \item 0.0500 \item 0.0200
  \end{enumerate}

  \item A 95\% confidence interval for $\mu$ with $n = 16$ uses which critical value $t^*$ ($df = 15$)?
  \begin{enumerate}[label=(\Alph*)]
    \item 1.960 \item 2.131 \item 1.753 \item 2.602 \item 2.947
  \end{enumerate}
\end{enumerate}

\begin{enumerate}[resume, leftmargin=1.5em]
  \item Which test evaluates whether categorical proportions are identical across 3 hospital networks?
  \begin{enumerate}[label=(\Alph*)]
    \item 1-sample $t$-test \item Chi-Square Test for Goodness-of-Fit \item Chi-Square Test for Homogeneity \item Chi-Square Test for Independence \item ANOVA
  \end{enumerate}

  \item What are the degrees of freedom for a regression $t$-test with $n = 28$?
  \begin{enumerate}[label=(\Alph*)]
    \item 28 \item 27 \item 26 \item 14 \item 2
  \end{enumerate}

  \item If the residuals of a model exhibit a curved U-shape, this indicates:
  \begin{enumerate}[label=(\Alph*)]
    \item A linear model is inappropriate \item The correlation is positive \item The slope is 0 \item The sample size is too large \item All residuals are zero
  \end{enumerate}

  \item A fair 6-sided die is rolled until a 4 appears. What is the probability it takes more than 3 rolls?
  \begin{enumerate}[label=(\Alph*)]
    \item $(1/6)^3$ \item $(5/6)^3 \approx 0.5787$ \item $1 - (5/6)^3$ \item $3(1/6)$ \item $(5/6)(1/6)^2$
  \end{enumerate}

  \item When constructing a confidence interval for $p$, which condition justifies assuming independence?
  \begin{enumerate}[label=(\Alph*)]
    \item $np \ge 10$ \item $n \le 0.10N$ \item $n \ge 30$ \item $r^2 > 0.8$ \item Double-blinding
  \end{enumerate}

  \item A 99\% confidence interval for $\mu_1 - \mu_2$ is $(1.2, 5.8)$. Based on this interval:
  \begin{enumerate}[label=(\Alph*)]
    \item $\mu_1$ is significantly greater than $\mu_2$ at $\alpha = 0.01$ because the entire interval is positive.
    \item $\mu_1 = \mu_2$.
    \item $\mu_2 > \mu_1$.
    \item No conclusion is possible.
    \item $p\text{-value} > 0.05$.
  \end{enumerate}

  \item The power of a hypothesis test can be increased by:
  \begin{enumerate}[label=(\Alph*)]
    \item Decreasing sample size $n$ \item Decreasing $\alpha$ \item Increasing sample size $n$ \item Making $H_0$ false \item Pooling variances
  \end{enumerate}

  \item In a completely randomized experiment with 60 subjects and 3 treatments, how many subjects per treatment if balanced?
  \begin{enumerate}[label=(\Alph*)]
    \item 10 \item 20 \item 30 \item 60 \item 180
  \end{enumerate}

  \item A student scores at the $84^{\text{th}}$ percentile on a normal exam with $\mu = 500, \sigma = 100$. The student's score is approximately:
  \begin{enumerate}[label=(\Alph*)]
    \item 584 \item 600 \item 650 \item 700 \item 750
  \end{enumerate}

  \item Which of the following is true about standard deviation $s_x$?
  \begin{enumerate}[label=(\Alph*)]
    \item It is resistant to outliers. \item It can be negative. \item It is measured in the same units as the original data. \item It equals the variance. \item It is unaffected by multiplication.
  \end{enumerate}

  \item The LSRL passes through which point?
  \begin{enumerate}[label=(\Alph*)]
    \item $(0, 0)$ \item $(\bar{x}, \bar{y})$ \item $(s_x, s_y)$ \item $(Q_1, Q_3)$ \item $(1, 1)$
  \end{enumerate}

  \item What type of bias occurs when voluntary callers participate in a radio poll?
  \begin{enumerate}[label=(\Alph*)]
    \item Nonresponse bias \item Voluntary response bias \item Stratification bias \item Wording bias \item Randomization bias
  \end{enumerate}

  \item Two independent events $A$ and $B$ have $P(A) = 0.3$ and $P(B) = 0.7$. What is $P(A \cap B)$?
  \begin{enumerate}[label=(\Alph*)]
    \item 1.00 \item 0.40 \item 0.21 \item 0.50 \item 0.00
  \end{enumerate}

  \item The $p$-value in a significance test represents:
  \begin{enumerate}[label=(\Alph*)]
    \item The probability that $H_0$ is true. \item The probability of obtaining a test statistic as extreme as observed, assuming $H_0$ is true. \item The margin of error. \item The probability of a Type II error. \item The power.
  \end{enumerate}

  \item In a $2 \times 3$ contingency table, degrees of freedom for Chi-Square test of independence is:
  \begin{enumerate}[label=(\Alph*)]
    \item 6 \item 5 \item 3 \item 2 \item 1
  \end{enumerate}
\end{enumerate}

\begin{enumerate}[resume, leftmargin=1.5em]
  \item If $s_x = 5, s_y = 10,$ and $r = -0.6$, what is the slope of the regression line?
  \begin{enumerate}[label=(\Alph*)]
    \item $-1.2$ \item $-0.3$ \item 1.2 \item $-0.6$ \item $-3.0$
  \end{enumerate}

  \item What is the minimum value for an expected count in a Chi-Square test?
  \begin{enumerate}[label=(\Alph*)]
    \item 1 \item 5 \item 10 \item 30 \item 50
  \end{enumerate}

  \item If $\bar{x} = 50, \mu_0 = 45, s = 10, n = 25$, what is the test statistic $t$?
  \begin{enumerate}[label=(\Alph*)]
    \item 1.0 \item 2.5 \item 5.0 \item 0.5 \item 2.0
  \end{enumerate}

  \item A double-blind experiment prevents which sources of bias?
  \begin{enumerate}[label=(\Alph*)]
    \item Both subject expectation bias and evaluator measurement bias \item Nonresponse bias \item Undercoverage \item Extrapolation \item Confounding from weather
  \end{enumerate}

  \item The sampling distribution of $\hat{p}$ has mean:
  \begin{enumerate}[label=(\Alph*)]
    \item $\bar{x}$ \item $\mu$ \item $p$ \item $np$ \item $\sqrt{p(1-p)}$
  \end{enumerate}

  \item If $z = -1.96$, the area to the left under the standard normal curve is:
  \begin{enumerate}[label=(\Alph*)]
    \item 0.05 \item 0.025 \item 0.95 \item 0.975 \item 0.50
  \end{enumerate}

  \item Which of the following conditions is required for regression inference?
  \begin{enumerate}[label=(\Alph*)]
    \item LINER \item BINS \item BITS \item CUSS \item PANIC
  \end{enumerate}

  \item An experiment that compares each participant's heart rate while resting versus after exercising is:
  \begin{enumerate}[label=(\Alph*)]
    \item Observational study \item Matched pairs design \item Completely randomized design \item Stratified survey \item Cluster trial
  \end{enumerate}

  \item If $P(A) = 0.8$, what is $P(A^c)$?
  \begin{enumerate}[label=(\Alph*)]
    \item 0.8 \item 0.2 \item 0.0 \item 1.0 \item 0.64
  \end{enumerate}

  \item The sum of residuals $\sum(y - \hat{y})$ for a least-squares regression line is:
  \begin{enumerate}[label=(\Alph*)]
    \item 1 \item $-1$ \item Always zero \item $r^2$ \item $s_y$
  \end{enumerate}

  \item A 90\% confidence interval for $\mu$ is $(12.4, 18.6)$. What is the margin of error?
  \begin{enumerate}[label=(\Alph*)]
    \item 6.2 \item 3.1 \item 15.5 \item 1.96 \item 0.90
  \end{enumerate}

  \item What type of variable is the number of cars passing an intersection in one hour?
  \begin{enumerate}[label=(\Alph*)]
    \item Categorical \item Continuous quantitative \item Discrete quantitative \item Qualitative \item Identifier
  \end{enumerate}

  \item If a significance test rejects $H_0$ at $\alpha = 0.05$, the $p$-value must be:
  \begin{enumerate}[label=(\Alph*)]
    \item $\le 0.05$ \item $> 0.05$ \item 0.05 exactly \item Negative \item 1.0
  \end{enumerate}

  \item When constructing a confidence interval, increasing the confidence level from 90\% to 99\%:
  \begin{enumerate}[label=(\Alph*)]
    \item Narrows the interval \item Widens the interval \item Has no effect on width \item Cuts sample size in half \item Decreases $z^*$
  \end{enumerate}

  \item The standard error of the slope $b_1$ measures:
  \begin{enumerate}[label=(\Alph*)]
    \item The true population slope $\beta$
    \item The estimated standard deviation of the sampling distribution of sample slopes $b_1$
    \item The correlation coefficient $r$
    \item The $y$-intercept
    \item The $r^2$ value
  \end{enumerate}
\end{enumerate}

\section{Section II: Free-Response Questions (Questions 1--6)}
\textbf{Time: 90 Minutes | 6 Questions | 50\% of Exam Score}

\begin{workbookbox}[Question 1: Data Exploration & Outlier Analysis (Short-Answer FRQ)]
A nutritionist records the sugar content (in grams per serving) for 14 brands of breakfast cereals:
\[ 4, \; 6, \; 7, \; 8, \; 9, \; 10, \; 11, \; 12, \; 12, \; 13, \; 14, \; 15, \; 18, \; 26 \]
(a) Determine the 5-number summary for this dataset.\\
(b) Calculate the outlier boundary fences using the $1.5 \times \text{IQR}$ rule and identify any confirmed outliers.\\
(c) Write a complete description of the distribution addressing center, spread, shape, and unusual features in context.
\end{workbookbox}

\begin{workbookbox}[Question 2: Experimental Design & Confounding (Short-Answer FRQ)]
A school district tests whether introducing standing desks increases student attentiveness. 40 high school classrooms are available for the study.
(a) Explain how you would implement a completely randomized design comparing standing desks against traditional desks.\\
(b) Describe a potential confounding variable and explain how a randomized block design (blocking by subject area, such as Math vs. History) would mitigate this confounding.
\end{workbookbox}

\begin{workbookbox}[Question 3: Probability & Random Variables (Short-Answer FRQ)]
A high-tech machine component has a 10\% failure rate within its first year. A company installs 8 independent components.
(a) What is the probability that exactly 2 components fail during the first year?\\
(b) What is the probability that at least 1 component fails during the first year?\\
(c) The replacement cost is \$150 per failed component plus a flat \$100 technician dispatch fee if any components fail. Calculate the expected total repair cost.
\end{workbookbox}

\begin{workbookbox}[Question 4: Inference for Proportions (Short-Answer FRQ)]
A university admissions office claims that at least 65\% of accepted students enroll. In an SRS of $n = 300$ accepted students, 180 enroll.
Conduct a complete hypothesis test at $\alpha = 0.05$ to evaluate whether the true enrollment proportion is significantly less than 65\%.
\end{workbookbox}

\begin{workbookbox}[Question 5: Inference for Means & Two-Sample Comparison (Short-Answer FRQ)]
An agricultural test compares the weight gain (in pounds) of cattle fed Diet 1 ($n_1 = 16, \bar{x}_1 = 85.4, s_1 = 6.2$) versus Diet 2 ($n_2 = 16, \bar{x}_2 = 79.1, s_2 = 5.8$). Boxplots of both samples show no outliers.
Construct and interpret a 95\% confidence interval for the true difference in mean weight gain ($\mu_1 - \mu_2$).
\end{workbookbox}

\begin{workbookbox}[Question 6: The Investigative Task (Extended Multi-Part FRQ)]
An educational research team models the relationship between daily study hours ($x$) and exam score ($y$) using regression: $\hat{y} = 42.0 + 7.5x$.
(a) Interpret the slope 7.5 in context.\\
(b) A student who studied 4 hours scored 78. Calculate and interpret the residual.\\
(c) The researchers discover that sleep duration ($w$, in hours) also impacts exam scores. When sleep is added, the model becomes $\hat{y} = 20.0 + 6.0x + 3.5w$. Explain how omitted variable bias affected the initial estimate of slope in part (a).\\
(d) Recommend an optimal study-and-sleep schedule subject to the constraint that study time plus sleep time cannot exceed 14 hours per day.
\end{workbookbox}

% =========================================================================
% PRACTICE EXAM 2 (2027 EXAM READY MASTERY EXAM)
% =========================================================================
\chapter{Practice Exam 2}

\section*{Exam Overview \& Instructions}
This second full-length practice exam provides comprehensive final review across all 9 curriculum units under strict testing conditions.
\begin{itemize}[leftmargin=1.2em]
  \item \textbf{Section I:} 40 Multiple-Choice Questions (90 Minutes | 50\% of Total Score).
  \item \textbf{Section II:} 6 Free-Response Questions (90 Minutes | 50\% of Total Score).
\end{itemize}

\section{Section I: Multiple-Choice Questions (Questions 1--40)}
\begin{enumerate}[leftmargin=1.5em]
  \item In a dataset of 100 values, the mean is 50 and standard deviation is 0. This implies:
  \begin{enumerate}[label=(\Alph*)]
    \item Half the values are 0 and half are 100. \item Every single value in the dataset equals 50. \item The variance is 50. \item The median is 0. \item The data is normal.
  \end{enumerate}

  \item If a distribution is skewed to the right, which is true?
  \begin{enumerate}[label=(\Alph*)]
    \item $\text{Mean} > \text{Median}$ \item $\text{Mean} < \text{Median}$ \item $\text{Mean} = \text{Median}$ \item Range is zero \item $s_x = 0$
  \end{enumerate}

  \item A scatterplot has $r = 0.0$. This means:
  \begin{enumerate}[label=(\Alph*)]
    \item There is no linear relationship between $x$ and $y$. \item There is no relationship of any kind. \item The data forms a vertical line. \item All points lie on the $y$-axis. \item $b_1 = 1$.
  \end{enumerate}

  \item In an experiment, the factor is:
  \begin{enumerate}[label=(\Alph*)]
    \item The explanatory variable manipulated by the researcher. \item The response variable. \item The subjects. \item The error rate. \item The randomizer.
  \end{enumerate}

  \item If $P(A) = 0.6, P(B) = 0.3,$ and $P(A \cap B) = 0.18$, events $A$ and $B$ are:
  \begin{enumerate}[label=(\Alph*)]
    \item Disjoint \item Independent \item Dependent \item Impossible \item Complements
  \end{enumerate}

  \item What is the expected value of a geometric random variable with $p = 0.25$?
  \begin{enumerate}[label=(\Alph*)]
    \item 0.25 \item 1.0 \item 2.0 \item 4.0 \item 16.0
  \end{enumerate}

  \item The sampling distribution of $\bar{x}$ will be approximately normal if:
  \begin{enumerate}[label=(\Alph*)]
    \item $n \ge 30$ by the Central Limit Theorem \item The population has $N \ge 30$ \item $p = 0.5$ \item $s_x = \sigma$ \item $\mu = 0$
  \end{enumerate}

  \item For a 95\% confidence interval for a proportion, the margin of error is:
  \begin{enumerate}[label=(\Alph*)]
    \item $1.96 \sqrt{\hat{p}(1-\hat{p})/n}$ \item $\hat{p} / n$ \item $1.96 / \sqrt{n}$ \item $\sqrt{p(1-p)}$ \item $2\sigma$
  \end{enumerate}

  \item When testing $H_0: \mu = 50$ vs. $H_a: \mu \neq 50$, $t = -2.00$ with $df = 20$. The two-sided $p$-value is:
  \begin{enumerate}[label=(\Alph*)]
    \item $2 \times P(t_{20} \le -2.00)$ \item $P(t_{20} \le -2.00)$ \item $1 - P(t_{20} \le -2.00)$ \item 0.05 \item 0.025
  \end{enumerate}

  \item Degrees of freedom for a Chi-Square goodness-of-fit test with 5 categories is:
  \begin{enumerate}[label=(\Alph*)]
    \item 5 \item 4 \item 25 \item 1 \item 10
  \end{enumerate}
\end{enumerate}

\begin{enumerate}[resume, leftmargin=1.5em]
  \item In a regression analysis, the slope $b_1 = -2.5$ and $\text{SE}(b_1) = 0.5$. The $t$-statistic for $H_0: \beta = 0$ is:
  \begin{enumerate}[label=(\Alph*)]
    \item $-5.0$ \item $-1.25$ \item 5.0 \item 0.20 \item $-0.20$
  \end{enumerate}

  \item If an observation has a negative residual, the actual data point:
  \begin{enumerate}[label=(\Alph*)]
    \item Lies below the regression line \item Lies above the line \item Lies on the line \item Is an outlier \item Has $z = 0$
  \end{enumerate}

  \item What happens to the power of a test when the sample size $n$ decreases?
  \begin{enumerate}[label=(\Alph*)]
    \item Increases \item Decreases \item Remains constant \item Becomes 1 \item Becomes 0
  \end{enumerate}

  \item A $z$-score of 0 indicates that the observation:
  \begin{enumerate}[label=(\Alph*)]
    \item Equals the mean \item Equals 0 \item Is an outlier \item Has a probability of 0 \item Is negative
  \end{enumerate}

  \item A survey draws an SRS of 100 students from a school of 1000. The 10\% condition is:
  \begin{enumerate}[label=(\Alph*)]
    \item Satisfied ($100 \le 0.10(1000) = 100$) \item Violated \item Not applicable \item Only for means \item Only for experiments
  \end{enumerate}

  \item Which of the following is an unbiased estimator of the population proportion $p$?
  \begin{enumerate}[label=(\Alph*)]
    \item Sample proportion $\hat{p}$ \item $\hat{p} + 0.05$ \item Sample mean $\bar{x}$ \item $s_x$ \item $z^*$
  \end{enumerate}

  \item The sum of squared residuals $\sum(y - \hat{y})^2$ is minimized by:
  \begin{enumerate}[label=(\Alph*)]
    \item The least-squares regression line \item The horizontal line through $\bar{y}$ \item The median line \item Any line with positive slope \item A polynomial curve
  \end{enumerate}

  \item If $P(A) = 0.4$, what is $P(A \cap A^c)$?
  \begin{enumerate}[label=(\Alph*)]
    \item 0.00 \item 0.40 \item 0.60 \item 1.00 \item 0.24
  \end{enumerate}

  \item For a 1-sample $t$-interval with $n = 30$, degrees of freedom is:
  \begin{enumerate}[label=(\Alph*)]
    \item 29 \item 30 \item 31 \item 28 \item 15
  \end{enumerate}

  \item In an experiment evaluating three pain relievers, patients are assigned by age group. Age is the:
  \begin{enumerate}[label=(\Alph*)]
    \item Blocking variable \item Response variable \item Treatment \item Placebo \item Dependent variable
  \end{enumerate}

  \item If $z^* = 2.576$, the corresponding confidence level is:
  \begin{enumerate}[label=(\Alph*)]
    \item 90\% \item 95\% \item 98\% \item 99\% \item 99.9\%
  \end{enumerate}

  \item A distribution of house prices has $\mu = \$450,000$ and $\text{Median} = \$320,000$. The shape is:
  \begin{enumerate}[label=(\Alph*)]
    \item Skewed right \item Skewed left \item Symmetric \item Uniform \item Normal
  \end{enumerate}

  \item What is the probability of flipping 4 heads in a row with a fair coin?
  \begin{enumerate}[label=(\Alph*)]
    \item 0.50 \item 0.25 \item 0.125 \item 0.0625 \item 0.03125
  \end{enumerate}

  \item A high leverage point in regression:
  \begin{enumerate}[label=(\Alph*)]
    \item Has an extreme $x$-value \item Always has a large residual \item Has $y = 0$ \item Is always an error \item Cannot affect the slope
  \end{enumerate}

  \item If $n = 64$ and $\sigma = 16$, the standard deviation of the sample mean $\bar{x}$ is:
  \begin{enumerate}[label=(\Alph*)]
    \item 2 \item 4 \item 8 \item 16 \item 0.25
  \end{enumerate}

  \item A significance test with $\alpha = 0.05$ produces $p\text{-value} = 0.082$. The decision is:
  \begin{enumerate}[label=(\Alph*)]
    \item Fail to reject $H_0$ \item Reject $H_0$ \item Accept $H_a$ \item Prove $H_0$ is true \item Increase $\alpha$
  \end{enumerate}

  \item In a Chi-Square test, the expected counts represent:
  \begin{enumerate}[label=(\Alph*)]
    \item Frequencies expected if the null hypothesis is true \item Actual counts collected \item Percentages of total \item Standard errors \item Degrees of freedom
  \end{enumerate}

  \item A matched pairs $t$-test is mathematically equivalent to:
  \begin{enumerate}[label=(\Alph*)]
    \item A 1-sample $t$-test on differences \item A 2-sample independent $t$-test \item A 1-prop $z$-test \item An ANOVA \item A Chi-Square test
  \end{enumerate}

  \item What percentage of observations fall within $\pm 1\sigma$ of the mean in a normal distribution?
  \begin{enumerate}[label=(\Alph*)]
    \item 50\% \item 68\% \item 95\% \item 99.7\% \item 100\%
  \end{enumerate}

  \item Which sampling method divides the population into clusters and samples all members of selected clusters?
  \begin{enumerate}[label=(\Alph*)]
    \item Cluster sampling \item Stratified sampling \item Simple random sampling \item Systematic sampling \item Convenience sampling
  \end{enumerate}

  \item If $P(A \mid B) = P(A)$, then $A$ and $B$ are:
  \begin{enumerate}[label=(\Alph*)]
    \item Independent \item Disjoint \item Mutually exclusive \item Dependent \item Equal
  \end{enumerate}

  \item Standard deviation of a binomial distribution is:
  \begin{enumerate}[label=(\Alph*)]
    \item $\sqrt{np(1-p)}$ \item $np$ \item $\sqrt{p(1-p)/n}$ \item $s / \sqrt{n}$ \item $\sigma^2$
  \end{enumerate}

  \item What does a residual plot with random scatter around zero indicate?
  \begin{enumerate}[label=(\Alph*)]
    \item A linear model is appropriate \item The model is curved \item High bias \item Outliers exist \item Extrapolation error
  \end{enumerate}

  \item A 95\% confidence interval for $p$ is $(0.40, 0.60)$. The point estimate $\hat{p}$ is:
  \begin{enumerate}[label=(\Alph*)]
    \item 0.50 \item 0.40 \item 0.60 \item 0.10 \item 0.20
  \end{enumerate}

  \item When testing $H_0: \mu = 0$ vs. $H_a: \mu > 0$, the test is:
  \begin{enumerate}[label=(\Alph*)]
    \item One-sided (right-tailed) \item One-sided (left-tailed) \item Two-sided \item Non-directional \item Chi-square
  \end{enumerate}

  \item A placebo is used in an experiment to:
  \begin{enumerate}[label=(\Alph*)]
    \item Control for the psychological effect of receiving a treatment \item Increase the sample size \item Ensure double-blinding \item Eliminate random assignment \item Guarantee significance
  \end{enumerate}

  \item If $n = 100$ and $p = 0.50$, is the Large Counts condition met?
  \begin{enumerate}[label=(\Alph*)]
    \item Yes ($np = 50 \ge 10$ and $n(1-p) = 50 \ge 10$) \item No \item Only for means \item Only if $N > 10,000$ \item Depends on $\sigma$
  \end{enumerate}

  \item In regression, $df$ for the slope $t$-test is:
  \begin{enumerate}[label=(\Alph*)]
    \item $n - 2$ \item $n - 1$ \item $n$ \item $2n$ \item $(r-1)(c-1)$
  \end{enumerate}

  \item The interquartile range (IQR) is:
  \begin{enumerate}[label=(\Alph*)]
    \item $Q_3 - Q_1$ \item $\text{Max} - \text{Min}$ \item $2s$ \item $\mu \pm \sigma$ \item $\bar{x} - M$
  \end{enumerate}

  \item An experimental design where each subject receives two treatments in random order is:
  \begin{enumerate}[label=(\Alph*)]
    \item Matched pairs design \item Completely randomized design \item Observational study \item Cluster trial \item Stratified survey
  \end{enumerate}
\end{enumerate}

\section{Section II: Free-Response Questions (Questions 1--6)}
\textbf{Time: 90 Minutes | 6 Questions | 50\% of Exam Score}

\begin{workbookbox}[Question 1: Comparative Data Analysis (Short-Answer FRQ)]
A regional hospital network compares patient recovery times (in days) between two rehabilitation protocols: Protocol Alpha ($n = 25$) and Protocol Beta ($n = 25$). 
The summary statistics are:
\begin{itemize}
  \item Protocol Alpha: $\text{Min} = 4$, $Q_1 = 7$, $\text{Median} = 10$, $Q_3 = 14$, $\text{Max} = 22$, $\bar{x} = 10.8$, $s = 4.5$.
  \item Protocol Beta: $\text{Min} = 6$, $Q_1 = 11$, $\text{Median} = 15$, $Q_3 = 18$, $\text{Max} = 29$, $\bar{x} = 15.2$, $s = 5.6$.
\end{itemize}
(a) Determine if there are any outliers in either protocol using the $1.5 \times \text{IQR}$ rule.\\
(b) Compare the two recovery time distributions addressing Center, Unusual features, Shape, and Spread in context.\\
(c) A patient requires a protocol where recovery takes no longer than 14 days. Based on the data, which protocol would you recommend? Justify your choice.
\end{workbookbox}

\begin{workbookbox}[Question 2: Non-Linear Modeling \& Transformations (Short-Answer FRQ)]
A marine biologist investigates the relationship between ocean depth ($x$, in hundreds of meters) and underwater light intensity ($y$, in lumens per square meter). The raw scatterplot displays a distinct non-linear decaying curve. 
A natural logarithm transformation is applied to the response variable, producing the least-squares regression equation:
\[ \widehat{\ln(y)} = 6.42 - 0.58x \]
The standard error of the slope is $s_{b_1} = 0.042$, and $r^2 = 0.941$ with $n = 22$ depth measurement stations.
(a) Predict the light intensity $y$ at a depth of 500 meters ($x = 5.0$).\\
(b) At a depth of 300 meters ($x = 3.0$), the observed light intensity was 120 lumens/$\text{m}^2$. Calculate and interpret the residual in the transformed scale.\\
(c) Interpret the coefficient of determination $r^2 = 0.941$ in context.
\end{workbookbox}

\begin{workbookbox}[Question 3: Random Variables \& Combinations (Short-Answer FRQ)]
A commercial express delivery service uses two independent routes to deliver packages between two distribution hubs.
Let $X$ represent the transit time (in hours) along Route 1, where $\mu_X = 4.2$ hours and $\sigma_X = 0.6$ hours.
Let $Y$ represent the transit time (in hours) along Route 2, where $\mu_Y = 5.1$ hours and $\sigma_Y = 0.8$ hours.
(a) A dispatcher dispatches one vehicle along Route 1 and one vehicle along Route 2. Find the mean and standard deviation of the total combined transit time $T = X + Y$.\\
(b) Find the mean and standard deviation of the difference in transit times $D = Y - X$.\\
(c) Assuming both transit times follow approximately normal distributions, calculate the probability that Route 2 takes at least 1.5 hours longer than Route 1.
\end{workbookbox}

\begin{workbookbox}[Question 4: Inference for Two Proportions (Short-Answer FRQ)]
A consumer advocacy organization surveys adult smartphone users to determine whether there is a difference in the proportion of satisfied users between Brand A and Brand B.
\begin{itemize}
  \item Brand A: In an SRS of 400 users, 328 report being satisfied.
  \item Brand B: In an SRS of 500 users, 375 report being satisfied.
\end{itemize}
(a) Construct a 95\% confidence interval for the true difference in the proportion of satisfied users ($p_A - p_B$).\\
(b) Based on your confidence interval, is there convincing statistical evidence that Brand A has a higher customer satisfaction rate than Brand B? Explain.
\end{workbookbox}

\begin{workbookbox}[Question 5: Inference for Means: Paired $t$-Test (Short-Answer FRQ)]
A corporate wellness director evaluates an 8-week mindfulness stress reduction workshop. Ten executives record their resting heart rate (beats per minute) immediately before starting the workshop and after completion:
\begin{center}
\begin{tabular}{|l|c|c|c|c|c|c|c|c|c|c|}
\hline
Executive & 1 & 2 & 3 & 4 & 5 & 6 & 7 & 8 & 9 & 10 \\ \hline
Before & 82 & 78 & 88 & 74 & 80 & 85 & 90 & 76 & 84 & 79 \\ \hline
After & 75 & 74 & 81 & 72 & 76 & 80 & 82 & 73 & 78 & 75 \\ \hline
Difference (B $-$ A) & 7 & 4 & 7 & 2 & 4 & 5 & 8 & 3 & 6 & 4 \\ \hline
\end{tabular}
\end{center}
The sample mean difference is $\bar{x}_d = 5.0\text{ bpm}$ and the sample standard deviation is $s_d = 1.886\text{ bpm}$. A normal probability plot of the differences shows a strong linear pattern with no outliers.
Carry out an appropriate significance test at the $\alpha = 0.01$ level to determine if the mindfulness workshop significantly reduces executive resting heart rate.
\end{workbookbox}

\begin{workbookbox}[Question 6: Investigative Task: Decision Analysis Under Uncertainty]
A biomedical startup is deciding whether to manufacture a rapid diagnostic test kit for a rare disease. 
The disease prevalence in the target population is $P(D) = 0.02$.
The diagnostic test has the following performance characteristics:
\begin{itemize}
  \item Sensitivity (True Positive Rate): $P(+ \mid D) = 0.95$
  \item Specificity (True Negative Rate): $P(- \mid D^c) = 0.98$
\end{itemize}
(a) Calculate the probability that a randomly chosen individual tests positive, $P(+)$.\\
(b) Calculate the Positive Predictive Value (PPV), which is the conditional probability that a person actually has the disease given that they received a positive test result, $P(D \mid +)$. Interpret this result in context.\\
(c) The medical clinic must adopt a treatment protocol. If a patient is treated when they actually have the disease, the net health benefit is $+$10,000. If an uninfected person is treated unnecessarily (false positive), adverse drug side effects incur a net cost of $-$2,500. If an infected person is left untreated (false negative), severe complications incur a cost of $-$50,000. Correctly withholding treatment from an uninfected person has a net payoff of $0.
Using your probabilities from parts (a) and (b), calculate the expected monetary/health payoff of adopting the rapid diagnostic test protocol for 1,000 tested individuals versus treating no one.
\end{workbookbox}